\section{Гидродинамика}
\subsection{Гидростатика}
%%%%%%%%%%%%%%%%%%%%%%%%%%%%%%%%%%%%%%%%%%%%%%%%%%%%
%%%%%%%%%%%%%%%%%%%%%%%%%%%%%%%%%%%%%%%%%%%%%%%%%%%%
 \z[{\begin{wrapfigure}[7]{r}{0.3\textwidth}%
		\vspace{-12pt}%
		\begin{flushright}%
 			\includegraphics[width=0.3\textwidth]{9.1.1}%
		\end{flushright}%
	\end{wrapfigure}%
	}]%
	{Шар перекрывает отверстие радиуса $r$ в плоской стенке, разделяющей жидкости, давление которых $3P$ и $P$. С какой силой $F$ прижимается шар к отверстию?}
%
{$F = 2\pi r^2 P$}
%%%%%%%%%%%%%%%%%%%%%%%%%%%%%%%%%%%%%%%%%%%%%%%%%%%%
\z{В полусферический колокол, края которого плотно прилегают к поверхности стола, наливают через отверстие вверху жидкость. Когда жидкость доходит до отверстия, она приподнимает колокол и начинает из-под него течь. Найдите массу $m$ колокола, если его внутренний радиус равен $R$, а плотность жидкости $\rho$.}
%
{$m = \frac{\pi R^3 \rho}{3}$}
%%%%%%%%%%%%%%%%%%%%%%%%%%%%%%%%%%%%%%%%%%%%%%%%%%%% 
 \z[{\begin{wrapfigure}[7]{r}{0.4\textwidth}%
		\vspace{-15pt}%
		\begin{flushright}%
 			\includegraphics[width=0.4\textwidth]{9.1.3}%
		\end{flushright}%
	\end{wrapfigure}%
	}]%
	{В вертикально расположенных цилиндрах, площади сечений которых $S_1$ и $S_2$, находятся два невесомых поршня, соединённых невесомой пружиной жёсткостью $k$. Пространство между поршнями заполнено водой. Нижний поршень (площадью $S_2$) в начальном состоянии поддерживается так, что пружина не напряжена, её длина при этом равна $l_0$ (рис. а). Затем поршень площадью $S_2$ отпускают, и оба поршня опускаются (рис. б). На какое расстояние $x$ сместится поршень площадью $S_1$? Изобразите графически зависимость $x(k)$ от жёсткости $k$ пружины. Оба цилиндра сообщаются с атмосферой.}
%
{$x = \frac{S_1 S_2}{(S_1 - S_2)^2} \cdot \frac{\rho g S_2 l_0}{k - \frac{\rho g S_1 S_2}{S_1 - S_2}}$ при $k \geq \frac{\rho g S_1^2 S_2}{(S_1 - S_2)^2};$ ~
$x = l_0$ при $k < \frac{\rho g S_1^2 S_2}{(S_1 - S_2)^2}$.}
%%%%%%%%%%%%%%%%%%%%%%%%%%%%%%%%%%%%%%%%%%%%%%%%%%%% 
 \z[{\begin{wrapfigure}[5]{r}{0.3\textwidth}%
		\vspace{-20pt}%
		\begin{flushright}%
 			\includegraphics[width=0.3\textwidth]{9.1.4}%
		\end{flushright}%
	\end{wrapfigure}%
	}]%
	{Закрытый цилиндр радиуса $R$, заполненный на три четверти своего объема жидкостью плотности $\rho$, вращается в невесомости вместе с жидкостью с угловой скоростью $\omega$ вокруг своей оси. Как меняется давление $p$ в жидкости в зависимости от расстояния $r$ до стенок цилиндра?}
%
{$P(x) = \frac{\rho \omega^2}{2} \left((R - x)^2 - \frac{R^4}{4}\right)$}
%%%%%%%%%%%%%%%%%%%%%%%%%%%%%%%%%%%%%%%%%%%%%%%%%%%%	
 \z[{\begin{wrapfigure}[6]{r}{0.3\textwidth}%
		\vspace{+5pt}%
		\begin{flushright}%
 			\includegraphics[width=0.3\textwidth]{9.1.5}%
		\end{flushright}%
	\end{wrapfigure}%
	}]%
	{В сосуде, дно которого образует угол $\alpha$ с горизонтом, стоит куб с ребром $a$, сделанный из материала плотности $\rho$. Верхнее ребро куба находится на глубине $h$. Жидкость под основание куба не подтекает. Атмосферное давление $P_0$, плотность жидкости $\rho_0$. Найдите силу $N$, с которой куб действует на дно сосуда.}
%
{$F_{\parallel} = a^3 (\rho - \rho_0) g \sin\alpha;$ ~
$F_\perp = a^3 \rho_0 g \left(\frac{\rho}{\rho_0} \cos\alpha + \frac{1}{2} \sin\alpha + \frac{h}{a}\right) + P a^2$}